\documentclass[10pt,a4paper,twocolumn]{article}

\usepackage[margin=20mm]{geometry}
\usepackage[T1]{fontenc}
\usepackage{lmodern}
\usepackage{microtype}
\usepackage{amsmath}
\usepackage{graphicx}
\usepackage{booktabs}
\usepackage{xcolor}
\usepackage[hidelinks]{hyperref}
\graphicspath{{figures/}}

% Keep this source self-contained for the Codex editor's PDF preview.
% Replace placeholders with evidence-backed prose as the report develops.
\newcommand{\todo}[1]{\textcolor{blue!60!black}{[TODO: #1]}}
\newcommand{\projecttitle}{Quantized EfficientNet-Lite0 Acceleration on RV64IMC}
\newcommand{\reportauthors}{Mateusz Kazimierczak, X1150017}

\title{\projecttitle\\\large Midterm Report}
\author{\reportauthors}

\begin{document}
\maketitle

\begin{abstract}
\todo{Summarize the workload, measured bottlenecks, proposed ISA extensions,
verification approach, and principal ISA-level results. Write this last.}
\end{abstract}

\section{Introduction and Objectives}
\label{sec:introduction}
\todo{Introduce the project motivation and midterm scope: custom instruction
semantics, Spike integration, functional verification, and instruction profiling.
State the questions addressed and the contributions demonstrated so far.}

\section{Workload and Experimental Setup}
\label{sec:setup}
\subsection{Model and Reference Implementation}
\todo{Describe EfficientNet-Lite0, model provenance, W8A8 quantization, dominant
operators, the frozen LiteRT Micro reference, and deterministic test inputs.
Define the extent of correctness validation and its limitations.}

\subsection{Configurations and Measurement Method}
\todo{Define the LiteRT baseline A0, custom compiler B1.1, final software B1.7 (B1),
and each custom ISA configuration included in this report. Record tool versions,
compiler flags, profiling boundaries, and repeatability checks. Explain that
Spike instruction counts do not establish CV-Wally cycle speedup.}

\section{Baseline Characterization}
\label{sec:baseline}
\todo{Present dynamic instruction count and mix, memory instructions, MACs per
retired instruction, operator contributions, and code size. Identify at least
three measured bottlenecks and connect each to an optimization opportunity.}

% Add an operator contribution plot or instruction-mix figure here.
% \begin{figure}[t]
%   \centering
%   \includegraphics[width=\columnwidth]{baseline-profile.pdf}
%   \caption{Baseline profile over the stated measurement interval.}
%   \label{fig:baseline-profile}
% \end{figure}

\section{ISA Design and Software Integration}
\label{sec:design}
\subsection{Packed Integer MAC}
\todo{Specify Xqmac8 encoding, signed lane interpretation, accumulator width,
overflow behavior, register effects, and legality constraints. Explain the Spike
implementation and kernel mapping using the measured bottlenecks.}

\subsection{Exact Requantization}
\todo{Specify Xrequant rounding, shifts, saturation, parameter representation,
and exceptional cases. Explain equivalence to the integer software reference
and the incremental change relative to the packed-MAC configuration.}

\subsection{Additional Extensions and Tradeoffs}
\todo{Choose which later extensions belong in the midterm scope. For each,
explain the residual bottleneck, semantics, software integration, and tradeoffs.
Remove this subsection if it is outside the selected scope.}

\section{Verification and Quantitative Evaluation}
\label{sec:evaluation}
\subsection{Functional Correctness}
\todo{Report directed and randomized instruction tests, operator or kernel
replay, reference comparisons, and ISA audits. Identify tested inputs and avoid
generalizing beyond the validation coverage.}

\subsection{Instruction-Level Results}
\todo{Populate Table~\ref{tab:results} with comparable measurements. Include a
microbenchmark, two major model operators, and full inference or a complete
block. Report custom instruction utilization, load/store instructions per MAC,
and code-size changes where applicable.}

\begin{table}[t]
\centering
\caption{Full-inference comparison. Dashes mark measurements to be supplied;
instruction reduction is relative to A0 over identical profiling boundaries.}
\label{tab:results}
\small
\begin{tabular}{lrrr}
\toprule
Config. & Instructions & Reduction & Code bytes \\
\midrule
A0 & --- & --- & --- \\
B1.1 & --- & --- & --- \\
B1 & --- & --- & --- \\
X1 & --- & --- & --- \\
X2 & --- & --- & --- \\
\bottomrule
\end{tabular}
\end{table}

\subsection{Discussion and Limitations}
\todo{Explain which bottlenecks changed, whether the applicable quantitative
target was met, and what remains limiting. Separate instruction reductions from
cycle, timing, and area claims; include hardware results only with supporting
measurements and clearly defined configurations.}

\section{Remaining Work and Conclusion}
\label{sec:conclusion}
\todo{Summarize the supported findings. State remaining CV-Wally integration,
cycle measurement, synthesis, or evaluation work based on the current project
status, with milestones and risks.}

\appendix
\section{Reproducibility Record}
\label{sec:reproduction}
\todo{Record the source commit and relevant local changes, model checksum,
configuration names, commands, and paths to raw results. Distinguish archived
measurements from any newly reproduced runs.}

% Use \bibitem entries here and \cite{key} in the text as sources are added.
% Inline references keep the editor preview independent of external files.
\begin{thebibliography}{9}
\bibitem{project}
Project assignment and clarifications, repository files
\texttt{instructions.md} and \texttt{clarifications.md}.
\end{thebibliography}

\end{document}
